\chapter{Methods and Research Design}

\hspace{\parindent} This chapter describes how the study will develop, train, and evaluate a Multi-Agent Reinforcement Learning (MARL) model for dynamic bus scheduling under non-ideal conditions. The aim is a reproducible pipeline that runs from a calibrated stochastic microsimulation of the EDSA Carousel corridor through to measurable performance results under both ideal and non-ideal operating conditions.

\section{Research Approach}

This study takes a \textbf{quantitative, simulation-based experimental} approach. Whether a control policy improves a transit system is a testable hypothesis, and answering it requires repeated measurement under controlled variation of inputs. Deploying an experimental dispatching policy directly on the operational EDSA Carousel fleet would carry unacceptable safety, financial, and service-continuity risks, so microsimulation serves as the experimental substrate. Two limitations are acknowledged upfront: results depend on simulation fidelity, addressed here through empirical calibration against real corridor data; and the modeled disturbances do not cover every real-world failure mode, addressed transparently in Scope and Limitations.

\section{Research Design}

The study uses a \textbf{controlled comparative experimental design} with two between-condition factors and three response variables. The factors are control strategy (\{No Control, Forward Headway, Even Headway, MARL\}) and operating condition (\{ideal, non-ideal\}). The response variables are mean passenger waiting time, mean total travel time, and headway coefficient of variation. The non-ideal condition is not a single setting but a range of disturbance intensities, so controller performance is characterized as a function of disturbance magnitude rather than evaluated at a single point.

Each combination is replicated $N \geq 30$ times using independent Monte Carlo seeds, so the Central Limit Theorem applies to the run means \cite{Wasserman2004}. The unit of analysis is one full simulated operating day on the defined EDSA Carousel sub-corridor. Random seeds and disturbance realizations are matched across control strategies within each disturbance level, so any performance gap between two controllers on the same run reflects the controller rather than the specific disturbance realization \cite{Henderson2018Matters,Agarwal2021Precipice}. Sampling is exhaustive within the simulation, so inclusion is defined at the disturbance level: only the modeled disturbances enter the non-ideal condition, and all others stay at their ideal values.

\section{Methodology}

\subsection{Simulation Environment}

The study uses a two-stage simulation architecture. \textbf{SUMO} (Simulation of Urban Mobility)~\cite{Lopez2018SUMO} serves as the \textbf{calibration layer}, providing a microscopic representation of the EDSA Carousel corridor from which empirical parameters are extracted. A \textbf{custom Python environment built on the PettingZoo Agent Environment Cycle (AEC) API}~\cite{Terry2021PettingZoo} serves as the \textbf{training and Monte Carlo evaluation layer}, where MARL agents interact with a lightweight bus dynamics model parameterized by the distributions extracted during SUMO calibration. SUMO is not used as the training environment itself for two reasons, its microscopic simulation of surrounding non-bus traffic is unnecessary for the agents' reward, which depends only on bus state, and stepping that traffic at every control event across the full Monte Carlo budget would impose a prohibitive computational cost. Hence SUMO is used once for calibration and the repeated training and evaluation runs are carried out in the lightweight Python environment.

PettingZoo's AEC API is chosen over three alternatives. First, the single agent Gymnasium interface~\cite{Towers2024Gymnasium} cannot typically represent multiple agents acting on independent events. Using multi-agent control into a single-agent API requires either treating the entire fleet as one composite agent which breaks parameter sharing or manual agent coordination outside the API which removes the benefit of using a standard interface at all. Second, PettingZoo's parallel API assumes all agents step simultaneously~\cite{Terry2021PettingZoo}, which does not match asynchronous bus control where only the bus arriving at a control stop has a decision to make. Third, a fully custom implementation without any standard API was considered but rejected because it would forfeit compatibility with established algorithm libraries and complicate reproducibility for follow-on work. AEC was designed specifically for the asynchronous case, cycling through one agent at a time rather than stepping all agents in parallel~\cite{Terry2021PettingZoo}. This matches the formulation used by Wang and Sun \cite{Wangsun} and Rodriguez et al.~\cite{Rodriguez2023Cooperative}, both of which implemented the same asynchronous method but in a fully custom event driven simulators. Adopting the AEC API preserves their event-driven dynamics while providing a recognized interface and compatibility with standard parameter sharing actor-critic libraries.

SUMO is used for the calibration layer because it integrates with Python via the Traffic Control Interface (TraCI) and supports the GEH and RMSE validation protocols described later in this section. Ollero et al.~\cite{Ollero2024EDSA} previously developed an AIMSUN microsimulation of the EDSA Busway corridor but this study implements an equivalent network in SUMO to take advantage of its open-source TraCI interface, while adopting the same calibration targets.

\subsubsection{Control-Stop Selection}

The study operates on a sub-corridor of the EDSA Carousel route. Within the chosen sub-corridor, only a subset of stops is designated as \textit{control stops} where MARL agents are permitted to act, the remainder are non-control stops that contribute stochastic dwell and delay behavior without agent intervention. Control stops are selected according to four criteria, following the design of Rodriguez et al.~\cite{Rodriguez2023Cooperative} and the minimal-intervention principle of Wang and Sun~\cite{Wangsun}:

\begin{enumerate}
    \item \textbf{Origin terminal.} The route origin is always a control stop, since it is the natural point at which initial headways are regularized before buses enter the corridor.
    \item \textbf{Onset of high-demand segments.} The first stop of each high demand segment is designated a control stop, so that the agent can act before demand-driven irregularity propagates downstream.
    \item \textbf{Low through-volume preference.} Stops with substantial through-passenger volume (many onboard riders not alighting) are \textit{avoided} as control stops, so that holding actions do not penalize large numbers of seated passengers already in transit.
    \item \textbf{No adjacent control hubs.} Where two high-demand hubs are adjacent, only the upstream hub is designated, so that holding penalties do not compound across two consecutive control actions.
\end{enumerate}

These criteria are data-driven but data-agnostic: they specify \textit{how} stops are chosen from demand and through-volume statistics, without presupposing \textit{which} stops the eventual dataset will yield. Once the operational data are processed (Section~\ref{subsec:data-pipeline}) and the per-stop demand and through volume profiles are known, the criteria are applied mechanically to produce the control-stop list for the study sub-corridor.

These two reasons are worth setting out in full. A pure-SUMO training loop, in which TraCI is queried at every control event during every training episode, was considered but not adopted. The first reason is \textbf{scope}: MARL training only needs the bus state variables that feed control decisions which includes position, headway, onboard load, and waiting queue. A full microscopic simulation of surrounding non-bus vehicles is useful for calibration realism but does not enter the reward function, which depends only on passenger-centric quantities derived from bus state. The second is \textbf{computational cost}, a full microscopic simulation steps thousands of non-bus vehicles every simulated second, and routing every training episode through TraCI introduces process boundary overhead on each control event. Across the evaluation parameters of this study which are four controllers, five disturbance levels, $N \geq 30$ Monte Carlo runs per cell, plus an ablation series, this overhead compounds into thousands of unnecessary CPU-hours. A pure-Python alternative without SUMO calibration was also considered but rejected: it would lack the empirically validated grounding for per-segment travel-time distributions that the GEH-and-RMSE-calibrated SUMO model provides. SUMOis used for empirically calibrated microscopic realism during the one-time calibration phase, Python for fast event-driven Monte Carlo during training and evaluation while keeping the experiment tractable on accessible hardware, including a single cloud-notebook GPU runtime. Established MARL bus-scheduling work, including Rodriguez et al.~\cite{Rodriguez2023Cooperative} and Wang and Sun~\cite{Wangsun}, takes the same route for the same reasons, using event-based custom simulators with calibration parameters supplied externally from empirical data.

The hybrid pipeline therefore proceeds in two phases. First, SUMO is calibrated against empirical bus data, and the resulting per-segment travel time, demand, and dwell time distributions are saved. Second, the Python environment loads these distributions and drives a lightweight event-based bus simulator in which the MARL agents train and are evaluated. Because the Python environment is parameterized from distributions extracted from the calibrated SUMO model, its statistical behavior matches SUMO at the level of per-segment travel time, demand, and dwell time. The Python environment does not reproduce SUMO's microscopic interactions among individual non-bus vehicles, but these interactions enter the agent's reward function only through their aggregate effect on inter-stop bus travel time, which is captured by the calibrated distributions. The Python environment matches SUMO in the statistics that the reward function actually depends on, while leaving out the microscopic detail that it does not. This is the same argument made by Rodriguez et al.~\cite{Rodriguez2023Cooperative}, who trained in an event-based mesoscopic simulator parameterized by empirical corridor data.

\subsubsection{Environment Model Validation}

Before MARL training begins, the SUMO calibration layer is statistically verified to approximate real-world bus operation along the corridor. Two industry-standard calibration metrics are used: the GEH statistic for bus volume on the dedicated corridor, and the Root Mean Square Error (RMSE) for bus speed trajectories. The calibration is restricted to the bus corridor itself, since the agents' state and reward depend only on bus dynamics; surrounding mixed-traffic flows do not enter the Python environment.

\subsubsection{Bus Volume Validation via GEH Statistic}

The GEH statistic, introduced for traffic-modelling acceptance testing in the UK Design Manual for Roads and Bridges \cite{DMRBV12}, measures the error margin between simulated and observed hourly bus volumes on individual corridor segments:

\begin{equation} GEH = \sqrt{\frac{2(M - C)^2}{M + C}} \label{eq:geh} \end{equation}

where $M$ is the hourly bus volume from the simulation and $C$ is the observed hourly bus count from empirical operational data. The environment is accepted as statistically reliable when $GEH < 5$ for more than 85\% of corridor segments, following the FHWA Traffic Analysis Toolbox calibration criterion~\cite{FHWA2019TAT} as applied in recent SUMO calibration work by Patil et al.~\cite{Patil2025Conformal}.

\subsubsection{Speed Trajectory Calibration via RMSE}

Volume calibration alone does not guarantee that buses move realistically. RMSE evaluates how closely simulated bus speed trajectories match empirical observations:

\begin{equation} RMSE = \sqrt{\frac{\sum_{i=1}^{n}(x_{obs} - x_{model})^2}{n}} \label{eq:rmse} \end{equation}

where $x_{obs}$ is the observed bus speed and $x_{model}$ is the simulated bus speed \cite{Ollero2024EDSA}. To minimize RMSE, microscopic SUMO parameters --- principally maximum desired speed and speed acceptance for the bus vehicle class --- are iteratively tuned to fit empirical speed profiles \cite{Ollero2024EDSA}. Once both the GEH threshold and RMSE minimization are achieved, the environment is treated as a validated digital twin of the operational scope.

\subsubsection{Real-Time Data Incorporation}

In a real-world deployment, the agent's state vector at each control event would be drawn from existing transit-data infrastructure: forward and backward headways $(h_i, h_{i+1})$ from Automatic Vehicle Location (AVL) feeds; onboard and waiting passenger counts from Automatic Passenger Counters (APC) and Automatic Fare Collection (AFC) terminals; and environmental flags from public weather APIs and incident-management systems \cite{Wangsun}. Within this study, the Python environment generates synthetic equivalents of these signals at each control event by querying the analytical bus model and combining outputs from the stochastic generators described below. The MARL policy is therefore exposed during training to signal modalities that map directly onto deployment-grade real-time data sources.

\subsection{Operating Conditions}

The study evaluates controllers under two distinct operating conditions, used consistently throughout the experimental design. The symbols used to define these conditions ($\eta$, $\sigma_d$, $\sigma_s$, and the breakdown rate $\lambda$) are collected in Table~\ref{tab:notation}.

\textbf{Ideal conditions.} The corridor operates under the empirical baseline observed in the operational dataset, with no additional disturbance layered on top. Specifically: the weather disturbance generator is disabled ($\eta = 0$), the breakdown generator is disabled ($\lambda = 0$), and the demand and traffic-speed scaling factors are fixed at their baseline means ($\sigma_d = 0$, $\sigma_s = 0$), so every bus encounters the empirical mean travel time and demand profile for its segment and time of day. The Python environment still injects normal day-to-day variability through the empirically calibrated per-segment travel-time and demand distributions, so the ideal condition is \textit{stochastic but not disturbed} --- comparable to a normal weekday with no incidents, severe weather, or surge events. This serves as the control condition for Stage A evaluation (SO 3): it isolates each controller's behavior under nominal operation, before disturbance-handling is tested.

\textbf{Non-ideal conditions.} One or more of the four stochastic generators described in the next section is activated. The full non-ideal evaluation (Stage B, SO 4) sweeps the weather disturbance intensity $\eta$ across five values, with the breakdown generator active alongside throughout. Demand and traffic-speed scaling remain active under both conditions to preserve realistic operational variability; the two conditions differ in whether the heavy-tailed weather and discrete breakdown generators are enabled.

Throughout this chapter, \textit{ideal conditions} and \textit{non-ideal conditions} refer to these operating states of the simulated environment. \textit{Baseline controllers} refers separately to the three non-MARL control strategies (No Control, Forward Headway, Even Headway) against which the MARL policy is benchmarked. The two terms are kept distinct: a condition is a state of the world; a controller is a choice of algorithm.

\subsection{Data Processing}

The MARL framework is trained and evaluated on a single stream of input data pre-processed before training: corridor operational data used to calibrate SUMO and to derive baseline parameters for the stochastic generators. A second stream, the synthetic state observations generated by the Python environment at runtime, is produced on-the-fly during training and consumed once per control event.

\subsubsection{Required Datasets}

A single pre-processed dataset is required:

\begin{itemize}

\item \textbf{Corridor bus operational data.} A per-trip record of EDSA Carousel bus operation along the study sub-corridor, collected over a continuous observation window of at least two weeks. The required fields are GPS-tracked vehicle location, boarding and alighting events, passenger occupancy, operating speed, and \textit{dwell time} at each stop --- the dwell time being the interval a bus spends stationary at a stop serving passengers, measured from the moment the doors open for boarding and alighting to the moment they close and the bus is ready to depart, exclusive of any holding time subsequently imposed by the controller. These records yield the empirical distributions of bus cruising speed, inter-stop travel time, and demand under ideal operating conditions, used both to calibrate SUMO and to define the baseline operating point of the stochastic generators. A crowdsourced operational dataset of this form --- for example, the kind collected from the EDSA Busway via the SafeTravelPH mobile application --- provides a representative model of the required structure, and the methodology described here is sized to operate on a dataset of comparable structure and duration. Should bus-volume coverage prove insufficient for the GEH validation step --- that is, if observed hourly bus counts at the segment level do not span enough of the operating day to support a reliable comparison against simulated counts --- supplementary records may be requested from MMDA, DOTr, or the involved bus operators under the Freedom of Information process.

\end{itemize}

Severe-weather conditions are not estimated from operational data in this study but are injected methodically as a controlled experimental variable, with disturbance magnitudes anchored to validated literature values rather than to a corridor-specific severe-weather sample. The rationale is twofold: a short empirical observation window cannot reliably resolve the variance of a heavy-tailed severe-weather distribution at per-segment granularity, and a swept-disturbance design characterizes controller robustness across a range of conditions rather than at a single uncertain calibration point. Because the disturbance conditions are random by construction, they are not drawn from any single fixed dataset: each Monte Carlo run generates a fresh, independently seeded realization of demand, traffic-speed, weather, and breakdown disturbances, so the controllers are evaluated against many distinct disturbance datasets rather than one. The single empirical dataset is required only to fix the \textit{baseline} (ideal-condition) operating point against which these random disturbances are layered; the held-out validation split (Stage 3 below) further confirms that the calibrated baseline generalizes beyond the records used to tune it. The full disturbance parameterization is described in Section~\ref{subsec:stochastic-vars}.

\subsubsection{Data Processing Pipeline}
\label{subsec:data-pipeline}

Pre-processing proceeds in three stages.

\textit{Stage 1: Cleaning.} Trip records with missing GPS coordinates, missing timestamps, negative inter-stop times, or trips that fail integrity checks (for example, a later stop served before an earlier one) are dropped. Remaining records are normalized to a common time zone. Because the operational observation window is short (on the order of two weeks), the data cannot support separate weekend or holiday demand profiles --- a window of this length yields too few weekend days and typically no holidays to estimate distinct calendars reliably. Records are therefore filtered to a single consistent service-day type (regular weekdays) and binned by time-of-day, which is the calendar dimension that actually drives demand and travel-time variation within the modeled operating day. The restriction to weekday operation is recorded as a scope limitation.

\textit{Stage 2: Empirical distribution extraction.} For each (segment, time-of-day) bin, the cleaned operational data yield the mean $\mu$ and standard deviation $\sigma$ of inter-stop travel time under baseline operating conditions, producing the baseline coefficient of variation $CV_0 = \sigma / \mu$ used to anchor the disturbance generators. Boarding and alighting counts per (stop, time-of-day) bin yield the baseline demand profile. These statistics are serialized to disk as parameter files that the Python environment loads at initialization.

\textit{Stage 3: Train/validation split for calibration.} The calibration data are split chronologically into a calibration set, used to tune SUMO parameters until the GEH and RMSE targets are met, and a held-out validation set, used to confirm that the calibrated environment generalizes beyond the data used to tune it.

The MARL agents themselves are not trained on these raw datasets directly. They are trained on synthetic state vectors $s_{i,t}$ produced by the Python environment at each control event, which is the only point in the simulation where an agent acts. Each $s_{i,t}$ is assembled at runtime by querying the analytical bus model and is consumed once by the policy network before being discarded; there is no persistent training dataset of state-action transitions stored on disk other than the standard replay buffer used by the learning algorithm.

\subsection{Stochastic Traffic Variables}
\label{subsec:stochastic-vars}

Four stochastic generators inject variability into the Python environment. Generators (i) and (ii) follow the perturbation framework of Wang and Sun~\cite{Wangsun}; the weather generator's heavy-tailed lognormal formulation follows Patil et al.~\cite{Patil2025Conformal}; the breakdown generator follows the rescheduling formulation of Cao et al.~\cite{Cao2022Train}. Table~\ref{tab:notation} collects the symbols used across this section and the MARL formulation that follows, so that each variable is defined in one place before it is used.

\begin{table}[htbp]
\centering
\caption{Notation used in the stochastic generators and the MARL formulation.}
\label{tab:notation}
\begin{tabularx}{\textwidth}{l X l}
\hline
\textbf{Symbol} & \textbf{Definition} & \textbf{Domain / units} \\
\hline
$\eta$ & Weather disturbance intensity; directly sets the coefficient of variation of disturbed inter-stop travel time ($CV = \eta$). $\eta = 0$ disables the weather generator. & dimensionless \\
$\sigma_d$ & Standard deviation of the per-episode passenger-demand scaling factor. $\sigma_d = 0$ fixes demand at its baseline mean. & dimensionless \\
$\sigma_s$ & Standard deviation of the per-episode traffic-speed scaling factor. $\sigma_s = 0$ fixes cruising speed at its baseline mean. & dimensionless \\
$\lambda$ & Rate parameter of the Poisson process governing bus breakdowns. $\lambda = 0$ disables breakdowns. & events/hour \\
$\mu$ & Baseline (empirical) mean inter-stop travel time for a given segment and time-of-day bin. & seconds \\
$\sigma$ & Baseline standard deviation of inter-stop travel time. & seconds \\
$CV_0$ & Baseline coefficient of variation, $CV_0 = \sigma/\mu$, used to anchor the disturbance generators. & dimensionless \\
$\mu_{ln}, \sigma_{ln}$ & Location and shape parameters of the lognormal weather-delay distribution (method-of-moments). & --- \\
$h^-,\ \hat{h}^+$ & Forward headway (gap to bus ahead) and estimated backward headway (gap to bus behind). & seconds \\
$H_0$ & Desired (scheduled) headway. & seconds \\
$\alpha_{i,t}$ & Holding-strength action of agent $i$, $\alpha_{i,t}\in\Omega=\{0.0,0.1,0.2,0.3,0.4\}$, scaled by $\Delta T$. & dimensionless \\
$\Delta T$ & Maximum holding duration that $\alpha_{i,t}$ scales. & seconds \\
$\gamma,\ \beta$ & Discount factor and event-based discount rate ($e^{-\beta t}$ replaces $\gamma$ for random inter-event times). & dimensionless \\
$N$ & Number of bus agents (also used for the Monte Carlo run count, $N\geq 30$, where noted). & count \\
\hline
\end{tabularx}
\end{table}

\subsubsection{Random Seed Control}

All four generators draw from a single seeded pseudo-random number stream maintained by the Python NumPy library. At the start of each Monte Carlo run, a master seed is set in the simulation driver script, which initializes NumPy's BitGenerator. For each (disturbance level, control strategy) cell, the list of $N \geq 30$ seeds is fixed in advance and stored as a configuration file in the project repository. Within a disturbance level, the same seed list is reused across all four control strategies, so the demand trace, disturbance realization, and breakdown timing for run $k$ are identical regardless of which controller is active. This matching enables paired statistical comparisons, where any difference in performance metrics between two controllers on run $k$ is attributable to the controller rather than to a lucky or unlucky disturbance realization \cite{Henderson2018Matters}. Seeds across runs are independent, and the choice $N \geq 30$ is made so that the Central Limit Theorem applies to the run-mean of each response variable. For settings with heavy-tailed disturbances, bootstrap-based confidence intervals are reported rather than relying solely on the normal approximation at $N = 30$ \cite{Agarwal2021Precipice}.

\subsubsection{Passenger Demand}

The baseline empirical transit demand is perturbed each episode by a scaling factor sampled from $\mathcal{N}(1, \sigma_d^2)$, clipped to $[1, 3]$, following Wang and Sun~\cite{Wangsun}. The asymmetric clip focuses the test regime on demand surges rather than symmetric variation, since demand drops produce well-behaved, lightly loaded operating conditions that do not stress-test the controller. The upper bound of 3 corresponds to roughly a tripling of baseline boarding rates, which spans the range observed during major event let-outs and severe-weather mode shifts. Sampling occurs at the start of each simulation run, producing varied passenger demand profiles across episodes and exposing the MARL agents to non-stationary boarding patterns.

\subsubsection{Traffic Delays}

Mean cruising speed between stops is adjusted dynamically using a scaling factor drawn from $\mathcal{N}(1, \sigma_s^2)$, clipped to $[0.8, 1.2]$, representing typical daily congestion friction \cite{Wangsun}. This generator is kept distinct from the heavier-tailed weather perturbation below.

\subsubsection{Weather-Induced Anomalies}

Severe-weather disturbances are injected methodically as a controlled experimental variable rather than estimated from a fixed empirical severe-weather sample. Inter-stop travel time under disturbance is sampled from a Lognormal distribution, which captures the right-skewed, heavy-tailed delays characteristic of urban traffic under inclement weather \cite{Patil2025Conformal}. The lognormal is chosen over two alternatives: a Gaussian distribution, which is symmetric and cannot represent the heavy right tail observed empirically in severe-weather travel times; and Gamma or Weibull distributions, which can produce heavy tails but lack the direct empirical validation against urban traffic data that Patil et al.~\cite{Patil2025Conformal} provide for the lognormal form. The lognormal's right-skew matches the operational observation that severe delays are bounded below by free-flow travel time but unbounded above. The distribution's mean is the baseline inter-stop travel time $\mu$ extracted from the operational data, and its variability is set by a disturbance intensity parameter $\eta$, which directly specifies the coefficient of variation:

\begin{equation} CV = \eta \label{eq:cv_eta} \end{equation}

The lognormal shape and scale parameters follow the method-of-moments derivation:

\begin{equation} \sigma_{ln} = \sqrt{\ln(\eta^2 + 1)} \label{eq:sigma_ln} \end{equation}

\begin{equation} \mu_{ln} = \ln(\mu) - \frac{\sigma_{ln}^2}{2} \label{eq:mu_ln} \end{equation}

Travel time is drawn as $T \sim \text{LogNormal}(\mu_{ln}, \sigma_{ln})$. Patil et al.~\cite{Patil2025Conformal} validated this lognormal parameterization against INRIX freeway data via the Kolmogorov-Smirnov test, reporting a close fit at the highest variability level they tested ($KS = 0.036$, $p = 0.94$ at $CV = 1.0$).

The disturbance intensity $\eta$ is swept across five values, $\eta \in \{0.0, 0.3, 0.6, 1.0, 1.3\}$, chosen to span and slightly exceed the empirically validated CV range of Patil et al.~\cite{Patil2025Conformal}, who evaluate travel-time scenarios at $CV \in \{0.1, 0.3, 0.5, 0.7, 1.0\}$ and report that the lower CV values correspond to stable, near-free-flow conditions while CV values approaching 1.0 correspond to their highest-variability operating regime. The basis for each swept value is therefore as follows. $\eta = 0$ disables the generator and recovers the ideal condition. $\eta \in \{0.3, 0.6\}$ fall inside the range Patil et al.\ validate directly, representing moderate travel-time variability. $\eta = 1.0$ is set to the upper end of the CV range that Patil et al.\ validate against INRIX data via the KS test; it is the most variable condition for which their lognormal fit is empirically confirmed, and is adopted here as the nominal severe-disturbance operating point. $\eta = 1.3$ extends beyond the validated range as an extrapolated stress test, included to probe controller behavior under conditions more extreme than any for which the underlying distribution has been empirically fitted; results at $\eta = 1.3$ are therefore interpreted as a robustness probe rather than as a calibrated severe-weather scenario. We note explicitly that the mapping from $\eta$ to a specific named weather severity (e.g.\ ``moderate rain'' versus ``typhoon-grade'') is not established by Patil et al., whose CV values are not labeled by weather category; the $\eta$ sweep should therefore be read as a span of travel-time variability magnitudes rather than as calibrated meteorological intensities. Reporting controller performance across this sweep characterizes robustness as a function of disturbance magnitude, which is methodologically stronger than evaluating at a single point and avoids reliance on a small corridor-specific severe-weather sample for parameter estimation. Anchoring the swept range to a directly weather-and-travel-time-calibrated source is a known limitation that constrains how literally the highest $\eta$ values can be interpreted.

\subsubsection{Bus Breakdowns}

Unlike the continuous perturbations above, a bus breakdown is a \textit{discrete event} that removes a bus from active simulation \cite{Cao2022Train,Shi2022DistDRL}. Breakdowns are triggered at random times sampled from a Poisson process with a configurable malfunction rate $\lambda$ (Table~\ref{tab:notation}). When a breakdown occurs at bus $b_k$, $b_k$ is removed from the active agent set, leaving an enlarged headway between the bus immediately upstream of $b_k$ (the trailing bus $b_{k-1}$, which is now closing on the gap that $b_k$ vacated) and the bus immediately downstream of $b_k$ (the leading bus $b_{k+1}$). The passengers that $b_k$ would have served are now left to accumulate at the stops ahead of $b_{k-1}$, so the trailing bus $b_{k-1}$ inherits this residual demand as it advances. To prevent $b_{k-1}$ from being overwhelmed by the doubled demand --- which would slow its dwell times, deepen the gap to $b_{k+1}$, and trigger bunching with the bus behind it --- \textbf{the agent controlling $b_{k-1}$ executes a forced stop-skipping action} at a heavily loaded downstream stop: alighting is permitted but boarding is prohibited, so $b_{k-1}$ sheds onboard passengers without absorbing the full waiting queue and can close the gap to $b_{k+1}$ more quickly. Assigning the action explicitly to a still-active agent keeps responsibility attributable. The Poisson assumption is a tractable approximation for the single-day simulation horizon, over which vehicle age, maintenance schedule, and mileage since last service can be treated as constant covariates.

\subsection{MARL Architecture and MDP Formulation}

This subsection builds the formal learning framework in two steps. It first introduces the Bellman recursion that underlies value-based reinforcement learning, then specifies the Markov game tuple, the per-agent observation and action spaces, the reward-function design space, and the proposed learning algorithm. The justification for choosing a multi-agent formulation over a single-agent alternative is presented in Chapter~2.

\subsubsection{Bellman Foundation}

At the core of any reinforcement learning method is the dynamic programming principle established by Bellman~\cite{Bellman1957}. The agent's objective is not to minimize immediate reward but to maximize the expected sum of discounted future rewards. The state-action value function $Q^\pi(s,a)$, which denotes the expected return from taking action $a$ in state $s$ and following policy $\pi$ thereafter, satisfies the Bellman optimality equation:

\begin{equation}
Q^*(s, a) = \mathbb{E} \left[ r(s, a) + \gamma \max_{a'} Q^*(s', a') \mid s, a \right]
\label{eq:bellman}
\end{equation}

where $s'$ is the state reached after applying action $a$, $r(s,a)$ is the immediate reward, and $\gamma \in [0, 1)$ is the discount factor. This equation matters here for two reasons. First, it says that a long-horizon control policy can be learned by recursively backing up local one-step rewards, which is what makes it computationally feasible for an agent to evaluate the downstream effects of a holding or skipping decision without explicitly enumerating every future stop on the route. Second, in the asynchronous bus-control setting, the time between two consecutive control events is itself random. Following Bradtke and Duff~\cite{Bradtke1995}, the discount factor in Eq.~\eqref{eq:bellman} is replaced by an event-based discount $e^{-\beta t}$ that scales with the elapsed time between events. Each agent's learning target is therefore an event-adjusted version of Eq.~\eqref{eq:bellman}, not a fixed-tick MDP target.

\subsubsection{Markov Game Formulation}

Following Littman~\cite{Littman1994}, the bus fleet control problem is formulated as a multi-agent extension of a Markov Decision Process, that is, a Markov game, defined by the tuple $G = \langle N, S, A, P, R, \gamma \rangle$:

\begin{itemize}

\item $N$: number of bus agents.

\item $S$: joint state space.

\item $A$: joint action space.

\item $P$: state-transition function $P(s_{t+1}\mid s_t, a_t)$.

\item $R$: reward function.

\item $\gamma \in [0, 1)$: discount factor.

\end{itemize}

Because bus control is asynchronous and event-driven, agents act only upon arrival at a control stop. The formulation maintains a per-agent local transition $P_i : S_i \times A_i \rightarrow S_i$, with $S_i$ the local observation of agent $i$ \cite{Wangsun,Rodriguez2023Cooperative}.

\subsubsection{Double Deep Q-Network Foundation}
\label{subsec:ddqn-foundation}

The Bellman optimality equation in Eq.~\eqref{eq:bellman} defines the value $Q^*(s,a)$ that an ideal controller would assign to each state-action pair, but it cannot be solved exactly for a state space as large as the bus-control problem, where the observation is a continuous vector of headways, loads, and queues. Two successive approximations bridge this gap. The first, $Q$-learning~\cite{Bellman1957,Littman1994}, replaces the exact solution with an iterative update that nudges a stored estimate $Q(s,a)$ toward the one-step Bellman target $r + \gamma \max_{a'} Q(s',a')$ after each observed transition. The second, the Deep $Q$-Network (DQN)~\cite{Mnih2015DQN}, replaces the lookup table with a neural network $Q(s,a;\theta)$ of weights $\theta$, so that the value function generalizes across the continuous state vectors the bus environment produces rather than requiring a separate stored entry for every possible state. DQN trains $\theta$ by minimizing the temporal-difference loss between the network's prediction and a target computed from a periodically frozen copy of the network, with transitions drawn from an experience replay buffer to decorrelate consecutive samples.

DQN has a known structural flaw that is decisive for this study: it systematically \textit{overestimates} action values~\cite{vanHasselt2016DDQN}. The cause is the single $\max$ operator in the target. Because the same network both proposes the next action ($\arg\max_{a'}$) and evaluates how good that action is ($\max_{a'} Q$), any noise in the value estimates is selected for --- the $\max$ preferentially picks whichever action happens to have been overvalued by chance, and that inflated value is then propagated back into the target. In a deterministic, well-behaved environment this bias is mild, but the swept-disturbance regime that this study intentionally generates is the opposite: heavy-tailed weather delays and discrete breakdowns produce high-variance returns, and high variance is precisely what feeds the overestimation mechanism. A controller trained under these conditions with plain DQN would be expected to learn systematically inflated value estimates, with no guarantee that the inflation is uniform across actions --- which is what makes it dangerous, since it can change which action looks best.

The Double Deep $Q$-Network (DDQN)~\cite{vanHasselt2016DDQN} removes this bias by \textit{decoupling action selection from action evaluation}. The online network with weights $\theta$ is still used to select the next action, but a separate target network with weights $\theta^-$ is used to evaluate that action's value. The DDQN learning target is

\begin{equation}
y^{\text{DDQN}} = r(s,a) + \gamma\, Q\!\left(s',\, \arg\max_{a'} Q(s', a'; \theta);\; \theta^-\right),
\label{eq:ddqn_target}
\end{equation}

in contrast to the plain-DQN target $y^{\text{DQN}} = r(s,a) + \gamma\, \max_{a'} Q(s', a'; \theta^-)$, which uses a single network's $\max$ for both roles. The distinction is that in Eq.~\eqref{eq:ddqn_target} the action is chosen by $\theta$ but scored by $\theta^-$: for the overestimation to survive, both networks would have to overvalue the \textit{same} action simultaneously, which is far less likely than a single network overvaluing it alone. The network is then trained by minimizing the squared temporal-difference error between its prediction and this target,

\begin{equation}
\mathcal{L}(\theta) = \mathbb{E}_{(s,a,r,s') \sim \mathcal{D}} \left[ \left( y^{\text{DDQN}} - Q(s,a;\theta) \right)^2 \right],
\label{eq:ddqn_loss}
\end{equation}

where $\mathcal{D}$ is the experience replay buffer and the target weights $\theta^-$ are held fixed for a number of updates (or slowly tracked toward $\theta$ via a soft update) so that the regression target does not move at every step.

Three properties make DDQN the appropriate value-based learner for this study, and each maps onto a specific feature of the bus-control problem. First, the discrete action space developed in Section~\ref{subsec:state-space} below ($|A_i| = 10$ holding-and-skipping actions per control event) is exactly the setting value-based methods are built for: the $\arg\max$ in Eq.~\eqref{eq:ddqn_target} is a finite enumeration over ten options, not a search over a continuous space. Second, the double estimator directly counters the overestimation bias that the study's deliberately high-variance disturbance regime would otherwise amplify, as argued above. Third, DDQN inherits DQN's experience-replay and target-network machinery, which is what makes it stable enough to train under disturbance without the auxiliary stabilizers that actor-critic methods require; the consequences of this choice for training stability are taken up in Section~\ref{sec:challenges}. The event-based discount $e^{-\beta t}$ introduced in the Bellman Foundation replaces the fixed $\gamma$ in Eq.~\eqref{eq:ddqn_target} when the target is computed in the asynchronous setting, so the DDQN target inherits the same event-adjustment as the underlying Bellman recursion. The combination of DDQN with parameter sharing and the centralized-training, decentralized-execution scheme is specified in the Proposed Learning Algorithm subsection.

\subsubsection{Agents ($N$)}

The decision-making agents are the individual bus units $i \in N$ operating along the corridor. To avoid combinatorial blow-up and to keep the system invariant to fleet size, all agents share a single neural-network policy under a \textbf{parameter-sharing} architecture and act decentrally on their own local observations at deployment \cite{Wangsun,Rodriguez2023Cooperative,Christianos2021PS}. Parameter sharing means that one network's weights serve every bus; each bus computes its own action from its own observation, but learning updates from any bus improve the policy for all of them.

\subsubsection{State Space and Local Observations ($S_i$)}
\label{subsec:state-space}

The local observation $s_{i,t}$ at agent $i$ on arrival at a control stop is a vector composed of:

\begin{itemize}

\item \textbf{Spatial location:} index of the current control stop.

\item \textbf{System regularity:} forward headway $h_i$ and backward headway $h_{i+1}$, both available from AVL feeds in deployment \cite{Wangsun}.

\item \textbf{Passenger demand:} onboard passenger count and estimated number of passengers waiting at the stop, available from APC/AFC systems \cite{Wangsun}.

\item \textbf{Environmental flags:} encoded indicators for the current disturbance intensity and any active downstream incident or breakdown.

\end{itemize}

\subsubsection{Action Space ($A_i$)}

\label{subsec:action-space} The agent's action $a_{i,t} = \pi_\theta(s_{i,t})$ combines two discrete components, following the action-space design of Rodriguez et al.~\cite{Rodriguez2023Cooperative}:

\begin{itemize}

\item \textbf{Holding control:} a holding-strength parameter $\alpha_{i,t} \in \Omega = \{0.0, 0.1, 0.2, 0.3, 0.4\}$, multiplied by a maximum holding duration $\Delta T$ to yield seconds held; $\alpha_{i,t} = 0$ means no holding.

\item \textbf{Stop-skipping:} a binary decision to bypass the next stop, triggered when the bus is severely delayed and near capacity, allowing headway recovery. Skipping is disallowed at the origin terminal and is not permitted if the immediately preceding trip already skipped that stop, so no passenger is denied service twice in a row.

\end{itemize}

The full action set is the Cartesian product of these two components: $|A_i| = 5 \times 2 = 10$ discrete actions per control event.

A continuous holding parameter $\alpha \in [0, 1]$ was considered, following the formulation of Wang and Sun~\cite{Wangsun}, but rejected for three reasons. First, continuous actions require actor-critic learning algorithms, whose training instability compounds across the swept-disturbance evaluation budget and would jeopardize timeline feasibility for a five-person team. Second, Rodriguez et al.~\cite{Rodriguez2023Cooperative} demonstrated that a 5-bin discretization of $\alpha$ achieves combined holding-and-skipping control on a comparable corridor without measurable loss of performance versus continuous formulations. Third, operational realism: real driver compliance with second-level holding instructions is itself coarse \cite{Rodriguez2023Cooperative}, so continuous precision in $\alpha$ is not meaningful at deployment.

\subsubsection{Reward Function ($R$) and Objective}
\label{subsec:reward}

The infinite-horizon objective of agent $i$ is the expected cumulative event-discounted return

\begin{equation}
R_i = \mathbb{E}\!\left[\sum_{k=0}^{\infty} e^{-\beta t_k}\, r_{i, t+k}\right],
\end{equation}

where $t_k$ is the elapsed time to the $k$-th future control event, consistent with the event-based discount introduced earlier.

The design of the per-event reward $r_{i, t+k}$ is itself a research objective of this study (Specific Objective~2, Expected Output~2.1) rather than an inherited choice. The reward formulation must balance three competing operational priorities: (i) headway regularity, since bunching is the primary failure mode of bus corridors; (ii) passenger waiting time at served stops, which is the most direct passenger-experience metric; and (iii) a penalty term that disincentivizes degenerate skipping behavior, since an agent given an unconstrained skip action could otherwise drive headway variance toward zero by refusing to serve passengers. Existing MARL bus-control reward formulations span a range of trade-offs among these three priorities, from headway-coefficient-of-variation forms designed for holding-only action spaces \cite{Wangsun,Wang2023MultiObj} to passenger-time forms designed for combined holding-and-skipping action spaces \cite{Rodriguez2023Cooperative}. This study will derive a reward function appropriate to the hybrid holding-and-skipping action space and the swept-disturbance operating regime, with the final form, parameter values, and sensitivity analysis reported in the implementation phase.

\subsubsection{Proposed Learning Algorithm}

The fully discrete action space ($|A_i| = 10$) makes value-based learning a natural fit. This study adopts \textbf{Double Deep Q-Network (DDQN)~\cite{vanHasselt2016DDQN} with parameter sharing under centralized training and decentralized execution (CTDE)}. CTDE means the policy is trained with access to information from all agents --- necessary for stable learning in non-stationary multi-agent environments --- but each bus acts on its own local observation alone at deployment, so the trained system does not require a centralized real-time controller. The choice follows Rodriguez et al.~\cite{Rodriguez2023Cooperative}, who used the same DDQN-based framework to learn a cooperative holding-and-skipping policy on a comparable bus corridor.

DDQN is chosen over three categories of alternatives. \textit{Plain DQN}~\cite{Mnih2015DQN} was considered but rejected: it overestimates Q-values in noisy environments \cite{vanHasselt2016DDQN,Henderson2018Matters}, which the high-variance disturbance regime intentionally generated by this study would aggravate. DDQN's double estimator addresses this directly by decoupling action selection from action evaluation \cite{vanHasselt2016DDQN}. \textit{Actor-critic methods such as MAPPO and MADDPG} were considered for their compatibility with continuous action spaces, but are unnecessary given the discrete action-space decision above; they would also introduce training instability that compounds across the swept-disturbance evaluation budget. \textit{Distributional MARL methods such as Wang and Sun's IQNC-M} \cite{Wangsun} explicitly model the return distribution rather than its expectation and offer the best published robustness on this problem class, but their meta-learning training procedure is at the research frontier and adds implementation risk inappropriate for a thesis-scale project; IQNC-M is therefore positioned as an extension direction in Chapter 5 rather than the primary algorithm here.

Three properties motivate the DDQN-with-parameter-sharing choice. First, DDQN's double estimator reduces the overestimation bias that plain DQN exhibits in noisy environments \cite{vanHasselt2016DDQN}, which matters because the swept-disturbance regime intentionally generates high-variance returns. Second, parameter sharing across agents preserves fleet-size invariance \cite{Christianos2021PS}: all bus agents share a single set of policy parameters, so the trained policy generalizes across changes in fleet composition. Third, DDQN training is more forgiving than actor-critic alternatives, which matters given the size of the evaluation budget --- training instability would compound across the swept-disturbance sweep.

Two enhancements common in the recent bus-control literature will be tested during implementation: extending agent awareness to include neighboring-trip information in the observation and reward, as in Rodriguez et al.'s DDQN-HA variant \cite{Rodriguez2023Cooperative}, and credit assignment via inductive graph learning, as in Wang and Sun's CAAC framework \cite{Wangsun}. The final architecture and any enhancements will be selected based on training stability and validation-set passenger waiting time after a fixed training budget.

\subsubsection{Training and Execution Protocol}

The learning process follows the standard reinforcement-learning feedback loop, instantiated here under the asynchronous, event-driven bus-control setting and the CTDE training scheme described above. The loop, repeated at every control event during training, proceeds as follows:

\begin{enumerate}

\item \textbf{Observe.} When bus agent $i$ arrives at a control stop, the Python environment assembles its local observation $s_{i,t}$ (Section~\ref{subsec:state-space}) from the analytical bus model and the active stochastic generators.

\item \textbf{Act.} The shared policy network selects an action $a_{i,t}$ (holding strength and skip decision) from $s_{i,t}$ using an $\epsilon$-greedy rule over the estimated $Q$-values during training, and greedily ($\arg\max_a Q$) at evaluation.

\item \textbf{Transition and reward.} The environment applies $a_{i,t}$, advances the simulation to the next control event, and returns the event-discounted reward $r_{i,t}$ (Section~\ref{subsec:reward}) reflecting the action's effect on headway regularity and passenger waiting.

\item \textbf{Store.} The transition $(s_{i,t}, a_{i,t}, r_{i,t}, s_{i,t'})$, where $t'$ is the agent's next control event, is written to a shared experience replay buffer common to all agents under parameter sharing.

\item \textbf{Update.} Mini-batches sampled from the replay buffer update the shared online network by minimizing the DDQN temporal-difference loss against a slowly-updated target network; experience from every bus updates the single shared parameter set.

\end{enumerate}

This closed loop --- observation in, action out, reward and next state back in, parameters updated --- is the feedback mechanism that distinguishes the learned controller from the fixed-rule baselines, which compute a holding time directly from the current headways with no update step.

The protocol is split into a training phase and an execution (evaluation) phase, consistent with the CTDE scheme:

\begin{itemize}

\item \textbf{Centralized training phase.} The loop above is run for a fixed episode budget with the stochastic disturbance generators \textit{active}, so that the policy is exposed to the same demand surges, traffic-speed friction, heavy-tailed weather delays, and discrete breakdowns that it will face at evaluation. Training under disturbance is a deliberate robustness measure: a policy that never observes a breakdown or a heavy-tailed delay during learning has no basis for responding to one at deployment, and would be expected to behave as an out-of-distribution controller under non-ideal conditions. Exposing the policy to randomized disturbances during training is the domain-randomization strategy that Tang et al.~\cite{Tang2024Curriculum} apply to bus operations for exactly this reason, randomizing variables such as passenger demand to obtain a policy that is robust to real-world traffic uncertainty. Training uses whatever system-wide information stabilizes learning (e.g.\ a shared reward and a shared buffer), and the experience of all buses contributes to one shared network.

\item \textbf{Decentralized execution phase.} The trained network's weights are frozen. During Monte Carlo evaluation (Stages A and B), each bus acts independently and greedily on its own local observation $s_{i,t}$ alone, with no centralized controller and no further weight updates. The fixed, already-trained policy is then evaluated across both the ideal condition (Stage A) and the swept non-ideal condition (Stage B), so that its behavior is measured on disturbance realizations it did not see during training but which are drawn from the same generative process it was trained against.

\end{itemize}

Because heavy-tailed returns can destabilize value-based learning if injected naively \cite{Henderson2018Matters,Agarwal2021Precipice}, training under disturbance is paired with the stabilizers discussed in Section~\ref{sec:challenges}: DDQN's double estimator to suppress overestimation bias, reward clipping, and slow target-network updates. If training under the full disturbance set proves unstable at the highest intensities despite these measures, a curriculum that ramps disturbance intensity over the course of training, in the manner of Tang et al.~\cite{Tang2024Curriculum}, is retained as a fallback and is noted as an implementation-phase contingency.

\subsection{Baseline Controllers for Comparison}

The MARL policy is benchmarked against three baseline controllers: \textbf{No Control (NC)}, \textbf{Forward Headway (FH)}, and \textbf{Even Headway (EH)}. These three are the standard benchmark set in the MARL bus-control literature; both Wang and Sun \cite{Wangsun} and Rodriguez et al.~\cite{Rodriguez2023Cooperative} use exactly this trio (in Rodriguez et al., EH is the primary benchmark). Using the same baselines makes the results of this study directly comparable to published work.

Each baseline represents a different point on the spectrum of control sophistication, from no intervention at all to a globally regularizing rule:

\subsubsection{No Control (NC)}

No holding or skipping is applied. Each bus dwells only for the time required to serve boarding and alighting passengers, then departs immediately. NC serves as the \textit{lower bound} of the comparison: it isolates the natural behavior of the corridor without any regularization mechanism. If MARL or any other controller cannot beat NC, the controller is doing something actively harmful. NC also provides the reference point against which the severity of bus bunching is measured, since bunching is the failure mode that NC fails to prevent.

\subsubsection{Forward Headway (FH)}

The classical holding rule introduced by Daganzo~\cite{Daganzo2009}. When a bus arrives at a control stop, its holding time is determined by how far behind the desired departure headway it is:

\begin{equation}
d_{FH} = \max\!\left\{0,\, \bar{d} + g\,(H_0 - h^-)\right\}
\label{eq:fh}
\end{equation}

where $H_0$ is the desired (scheduled) headway, $h^-$ is the bus's forward headway (time gap to the bus ahead), $\bar{d}$ is the average slack delay at equilibrium, and $g > 0$ is a control gain parameter. The intuition is straightforward: if a bus is running closer to the bus ahead than the schedule calls for ($h^- < H_0$), it holds for an additional period proportional to the gap deficit; if it is already at or beyond the desired spacing ($h^- \geq H_0$), it does not hold. Parameters $\bar{d}$ and $g$ follow the values used in Daganzo's original formulation. FH represents \textit{local, single-agent reactive control}: each bus reacts only to information about the bus immediately ahead, with no awareness of the bus behind or of system-wide state.

\subsubsection{Even Headway (EH)}

A more globally aware rule that holds buses to equalize the forward and backward gaps simultaneously, following the formulation of Rodriguez et al.~\cite{Rodriguez2023Cooperative}:

\begin{equation}
T_{EH} = \max\!\left\{\min\!\left[\frac{\hat{h}^+ - h^-}{2},\; 0.4\,H_0\right],\, 0\right\}
\label{eq:eh}
\end{equation}

where $h^-$ is again the forward headway, $\hat{h}^+$ is the \textit{estimated} backward headway (the time gap to the bus behind), and the term $\frac{\hat{h}^+ - h^-}{2}$ is the holding time that would split the difference evenly between the two gaps. The hold is capped at $0.4\,H_0$ (40\% of the scheduled headway), the bound Rodriguez et al.\ used to prevent excessive intervention; this cap also matches the maximum value in the MARL discrete action set $\Omega = \{0, 0.1, 0.2, 0.3, 0.4\}$, so EH and MARL are constrained to the same maximum hold and the comparison is fair. EH represents \textit{globally-aware reactive control}: it uses information about both neighbors but follows a fixed rule rather than a learned policy.

\subsubsection{Why These Three}

The three controllers are deliberately chosen to span the comparison space. \textbf{NC tests whether intervention helps at all.} \textbf{FH tests whether MARL outperforms classical local reactive control.} \textbf{EH tests whether MARL outperforms classical global reactive control} --- and EH is the stronger of the two classical baselines, consistently reported as the best non-RL controller in the literature \cite{Wangsun,Rodriguez2023Cooperative}. The acceptance criterion for the MARL policy in Stage A (Section~\ref{subsec:evaluation}) is therefore stated relative to EH, not NC: matching or beating EH is the meaningful test, since beating NC only proves the controller is not actively harmful.

This three-way comparison also lets the study attribute MARL's advantage. If MARL beats NC but not FH or EH, then any benefit is from the act of intervening, not from learning. If MARL beats FH but not EH, the benefit is from global awareness, not from learning. Only if MARL beats EH does the benefit attribute to the learned policy itself --- to the agent's ability to choose holding strengths and skipping decisions that a fixed rule cannot. The same logic carries into Stage B (SO 4): the disturbance-resilience claim is meaningful only if MARL's advantage over EH \textit{grows} with $\eta$, since both controllers see the same disturbance realizations.

State-of-the-art MARL methods such as CAAC \cite{Wangsun} and IQNC-M \cite{Wangsun} were considered as additional baselines but excluded from the primary comparison set for two reasons. First, including them would shift the research question from "does MARL help on EDSA?" (the SO 3 / SO 4 question) to "which MARL variant is best on EDSA?" (a different question requiring different data). Second, implementing CAAC or IQNC-M to publication-quality fidelity inside the thesis timeline would consume budget better spent on the EDSA-specific calibration and the disturbance-sweep analysis. They are flagged as comparison targets for follow-on work in Chapter 5.

\subsection{Data Analysis Methods}

For each (control strategy, disturbance level) cell, $N \geq 30$ independent Monte Carlo runs are executed using matched random seeds across strategies. Three response variables are logged per run: mean passenger waiting time, mean total travel time, and headway coefficient of variation. The choice $N \geq 30$ is made so that the Central Limit Theorem applies to the run-mean of each response variable \cite{Wasserman2004}, supporting standard statistical inference.

Statistical significance of performance differences between controllers will be assessed using hypothesis tests appropriate to the experimental design, at a significance level of $\alpha = 0.05$. Because the same random seeds are reused across controllers within a disturbance level, comparisons between two controllers on the same disturbance level are naturally \textit{paired}: the only thing differing between two runs is the controller itself, not the disturbance realization. Paired-sample tests will therefore be used for within-disturbance-level comparisons; unpaired tests will be used where seed matching does not apply (for example, comparing the same controller across different disturbance levels). Multiple-comparison correction will be applied across the four controllers to control the family-wise error rate, since running multiple pairwise tests inflates the chance of spurious significance.

For non-ideal disturbance levels, performance \textit{degradation} is computed as the percentage change in each response variable relative to the same controller's $\eta = 0$ ideal-condition result. Confidence intervals on this degradation ratio will be computed by bootstrap resampling, since ratios of sample means are not guaranteed to follow a normal distribution and may be skewed under high-disturbance conditions \cite{EfronTibshirani1993}. The specific hypothesis tests, normality screens, multiple-comparison procedure, and bootstrap parameters will be selected during the implementation phase in consultation with the research supervisor, based on the observed distributions of the per-run response variables and following recent guidance on rigorous reinforcement-learning evaluation \cite{Henderson2018Matters,Agarwal2021Precipice}. Analysis is performed in Python using NumPy, SciPy, and pandas; raw run logs and analysis notebooks are version-controlled to support reproducibility.

\subsection{Evaluation Methods}
\label{subsec:evaluation}

Evaluation proceeds in two stages tied to Specific Objectives 3 and 4.

\textbf{Stage A: Ideal-condition evaluation (SO 3).} The stochastic generators for severe weather and breakdowns are disabled; passenger demand and traffic-speed scaling are fixed at their baseline means (equivalent to $\eta = 0$). The trained MARL policy and the three baseline controllers (No Control, Forward Headway, Even Headway) are each run for $N \geq 30$ Monte Carlo iterations with matched seeds. The acceptance criterion for the MARL policy is twofold: (i) mean passenger waiting time no worse than Even Headway (no statistically significant degradation relative to EH at $p < 0.05$, with multiple-comparison correction applied across the four controllers), and ideally a statistically significant improvement over EH; and (ii) a statistically significant reduction in headway coefficient of variation relative to No Control. The two-part criterion reflects that under the well-behaved ideal condition EH is already near-optimal, so parity with EH is an acceptable Stage A outcome; the stronger ``beats EH'' claim is the target but is not required for the policy to pass Stage A, since the principal contribution of this study is the disturbance-robustness behavior characterized in Stage B.

\textbf{Stage B: Non-ideal-condition evaluation (SO 4).} The trained MARL policy and the three baseline controllers are evaluated across the swept disturbance range $\eta \in \{0.3, 0.6, 1.0, 1.3\}$, with the breakdown generator activated alongside the weather generator at each level. For each $(\eta, \text{controller})$ cell, $N \geq 30$ Monte Carlo iterations are run with matched seeds across controllers. Reported outputs are: (a) mean waiting time, travel time, and headway CV as functions of $\eta$ for each controller, plotted as performance curves; (b) percentage degradation at each $\eta$ relative to the controller's own $\eta = 0$ ideal-condition result; (c) the disturbance level at which MARL's advantage over the best-performing baseline controller becomes statistically significant under the chosen multiple-comparison correction; and (d) the degradation slope, defined as the rate of waiting-time increase per unit increase in $\eta$, with smaller slopes indicating greater robustness. An ablation is also run with each disturbance enabled alone (demand-only, weather-only, breakdown-only) at $\eta = 1.0$ to attribute degradation to disturbance type. The MARL policy is reported as robust if its mean waiting time stays below the mean waiting time of the best-performing baseline controller across the full sweep, with statistical significance under the same correction.

\section{Methodological Challenges}
\label{sec:challenges}

Three challenges shape the experimental design.

\textbf{Q-value overestimation under heavy-tailed returns.} Value-based methods such as DQN are known to overestimate Q-values \cite{vanHasselt2016DDQN}, and this bias is amplified when returns are heavy-tailed, as they are under high-intensity disturbance \cite{Henderson2018Matters,Agarwal2021Precipice}. The DDQN architecture's double estimator directly addresses this by decoupling action selection from action evaluation \cite{vanHasselt2016DDQN}. Because disturbances are active during training in this study, reward clipping and target-network soft updates are used as additional stabilizers throughout training, not only at the highest disturbance intensities.

\textbf{Computational cost.} The swept-disturbance design requires $N \geq 30$ Monte Carlo runs per (controller, disturbance level) cell. With four controllers and five disturbance levels, the Stage B evaluation budget exceeds 600 runs, with an additional ablation series ($\sim$360 runs) on top. Two architectural decisions keep this tractable on accessible hardware. First, SUMO is restricted to one-time calibration; the 600-plus Monte Carlo runs execute in the lightweight Python environment, which avoids the per-tick microscopic-vehicle simulation that dominates SUMO's cost. Second, evaluation runs are CPU-bound and embarrassingly parallel across random seeds, so training is the only phase requiring sustained GPU time.

\textbf{Timeline feasibility.} The full evaluation suite of approximately 1{,}080 simulation runs (Stage A, Stage B, ablation series combined) requires on the order of single-digit hours of single-core CPU time, given event-driven simulation rates reported in comparable bus-control work \cite{Wangsun,Rodriguez2023Cooperative}; parallelizing across seeds on a multi-core cloud-notebook runtime fits the full sweep within a single working session. MARL training is sized to the 500--800 episode budget reported in Rodriguez et al.~\cite{Rodriguez2023Cooperative}, with each full training run on the order of several hours of T4-class GPU time. The total training budget across the implementation phase, including hyperparameter tuning and reward-function iteration (Specific Objective 2.1), is approximately 30 hours of GPU time. This fits comfortably within the monthly compute allocation of standard cloud-notebook subscriptions and does not require institutional HPC resources. The longest single-path component is SUMO calibration, which depends on iterative parameter tuning to meet GEH and RMSE targets. A tiered scope plan is maintained throughout the implementation phase: the minimum deliverable is Stage A (ideal-condition evaluation, SO 3) plus a three-level non-ideal sweep at $\eta \in \{0, 0.6, 1.0\}$; the full five-level sweep and ablation series are added if calibration and training proceed on schedule. This tiering ensures that a complete thesis is deliverable even under conservative assumptions about calibration time.